\section{Detailed Cross-Testbed and Model-Family Results}
\label{app:cross_testbed_full_tables}

This appendix preserves the full cross-testbed and model-family tables that support the concise Results summaries in \Cref{tab:testbed_comparison_summary,tab:model_family_comparison_summary}. These tables are retained for auditability but moved out of the main Results flow so the paper body can focus on the central claims and final graph slots.

\subsection{Full Cross-Testbed Results}
\label{app:full_cross_testbed_results}
\shee{Add some text here and link to the corresponding tables or remove these sub-section titles}

\begin{table*}[t]
\centering
\caption{Full cross-testbed performance comparison using scenario-aggregated Oracle-normalized efficiency (\%) and scenario champions from \texttt{scenario\_winner}. Numeric columns are means across all five scenarios and four allocators. \paperTwo, \paperSeven, and \paperTwelve use the full evaluation corpora aggregated across $s\in\{1.0,1.5,2.0\}$ (\texttt{cap\_type}=T) with five runs per setting; \paperEight uses the current original paper-config slice at 1K/1K/1R, $s=1.0$, \texttt{cap\_type}=T.}
% \devroop{fix the testbed column dimensions. its overflowing.} -- SOLVED: Converted the appendix table to fixed-width columns with tighter spacing and shorter testbed descriptors.
\label{tab:testbed_comparison_full}
\scriptsize
\setlength{\tabcolsep}{2.5pt}
\renewcommand{\arraystretch}{1.08}
\begin{tabular}{@{}>{\raggedright\arraybackslash}p{0.19\textwidth} >{\raggedright\arraybackslash}p{0.19\textwidth} r r r r r@{}}
\toprule
\textbf{Testbed} & \textbf{Algorithm} & \textbf{Avg Reward} & \textbf{Regret} & \textbf{Eff.} & \textbf{Gap} & \textbf{Wins} \\
\midrule
\multirow{6}{=}{\textbf{\paperTwo}\newline 15N/51E/8P\newline 4K/2K/5R\newline 4 allocs, $s\in\{1,1.5,2\}$}
& \texttt{ORACLE} & 0.3926 & 0.0 & --- & --- & --- \\
& \texttt{CPursuitNeuralUCB} & 0.2887 & 509.2 & 73.22 & 26.78 & 38/300 \\
& \texttt{GNeuralUCB} & 0.2886 & 515.9 & 73.19 & 26.81 & 84/300 \\
& \texttt{iCPursuitNeuralUCB} & 0.2933 & 494.5 & \textbf{74.45} & 25.55 & \textbf{95/300$^\star$} \\
& \texttt{EXPNeuralUCB} & 0.2810 & 583.2 & 71.25 & 28.75 & 83/300 \\
& \cellcolor{findingBlue}\textit{Takeaway} & \multicolumn{5}{>{\raggedright\arraybackslash}p{0.51\textwidth}@{}}{\cellcolor{findingBlue}All pursuit hybrids remain competitive; iCPursuitNeuralUCB has the highest average efficiency and most wins.} \\
\midrule
\multirow{6}{=}{\textbf{\paperSeven}\newline 50N/141E/15P\newline 50/50/5R\newline 4 allocs, $s\in\{1,1.5,2\}$}
& \texttt{ORACLE} & 8.9237 & 0.0 & --- & --- & --- \\
& \texttt{iCPursuitNeuralUCB} & 6.9793 & 42.8 & \textbf{78.03} & 21.97 & \textbf{245/300$^\star$} \\
& \texttt{EXPNeuralUCB} & 6.2227 & 272.3 & 69.57 & 30.43 & 16/300 \\
& \texttt{CPursuitNeuralUCB} & 6.3314 & 269.1 & 70.82 & 29.18 & 0/300 \\
& \texttt{GNeuralUCB} & 6.3314 & 269.1 & 70.82 & 29.18 & 39/300 \\
& \cellcolor{findingGreen}\textit{Takeaway} & \multicolumn{5}{>{\raggedright\arraybackslash}p{0.51\textwidth}@{}}{\cellcolor{findingGreen}iCPursuitNeuralUCB dominates the dense topology in both efficiency and configuration-level wins.} \\
\midrule
\multirow{6}{=}{\textbf{\paperTwelve}\newline 100N/426E/4P\newline 1.5K/500/5R\newline 4 allocs, $s\in\{1,1.5,2\}$}
& \texttt{ORACLE} & 0.8724 & 0.0 & --- & --- & --- \\
& \texttt{GNeuralUCB} & 0.3833 & 864.1 & 43.72 & 56.28 & 53/300 \\
& \texttt{EXPNeuralUCB} & 0.3730 & 892.2 & 42.54 & 57.46 & 91/300 \\
& \texttt{iCPursuitNeuralUCB} & 0.3869 & 858.3 & \textbf{44.14} & 55.86 & \textbf{97/300$^\star$} \\
& \texttt{CPursuitNeuralUCB} & 0.3842 & 861.2 & 43.81 & 56.19 & 59/300 \\
& \cellcolor{findingOrange}\textit{Takeaway} & \multicolumn{5}{>{\raggedright\arraybackslash}p{0.51\textwidth}@{}}{\cellcolor{findingOrange}The mid-scale fusion topology compresses all methods, but iCPursuitNeuralUCB still leads.} \\
\midrule
\multirow{6}{=}{\textbf{\paperEight}\newline 20N/19E/8P\newline 1K/1K/1R\newline 4 allocs, $s=1$}
& \texttt{ORACLE} & 300540.1887 & 0.0 & --- & --- & --- \\
& \texttt{iCPursuitNeuralUCB} & 210418.9742 & 60833.4 & \textbf{67.85} & 32.15 & 1/20 \\
& \texttt{EXPNeuralUCB} & 186656.1295 & 82175.2 & 61.93 & 38.07 & \textbf{10/20$^\star$} \\
& \texttt{GNeuralUCB} & 183038.7972 & 83361.0 & 61.65 & 38.35 & 4/20 \\
& \texttt{CPursuitNeuralUCB} & 182011.6324 & 84024.2 & 61.42 & 38.58 & 5/20 \\
& \cellcolor{findingBlue}\textit{Takeaway} & \multicolumn{5}{>{\raggedright\arraybackslash}p{0.51\textwidth}@{}}{\cellcolor{findingBlue}iCPursuitNeuralUCB leads average efficiency, while EXPNeuralUCB wins more individual configurations.} \\
\bottomrule
\end{tabular}

\vspace{0.2cm}
{\scriptsize \textit{Legend:} N=nodes, E=edges, P=paths, R=runs. Efficiency (\%) = model average reward divided by Oracle average reward. Gap = 100 - Efficiency. Wins denote experiment-level wins out of the available configurations: 300 for \paperTwo/\paperSeven/\paperTwelve and 20 for \paperEight. $^\star$ indicates the testbed-level experiment-winner count.}
\end{table*}

\subsection{Full Model-Family Results}
\label{app:full_model_family_results}

\begin{table*}[t]
\centering
\caption{Full model-family performance comparison. The top block reports representative algorithms from each bandit family on the internal evaluation corpus; the bottom block reports hybrid pursuit-neural results on four external testbeds. Best per-family/testbed performance is bolded.}
\label{tab:model_family_comparison_full}
\scriptsize
\setlength{\tabcolsep}{2.8pt}
\renewcommand{\arraystretch}{1.08}
\begin{tabular}{@{}>{\raggedright\arraybackslash}p{0.18\textwidth} >{\raggedright\arraybackslash}p{0.26\textwidth} r r r r@{}}
\toprule
\textbf{Scope} & \textbf{Algorithm} & \textbf{Avg Eff.} & \textbf{Gap} & \textbf{Floor} & \textbf{Wins} \\
\midrule
\multirow{6}{=}{\textbf{CMABs}\newline internal corpus}
& \texttt{CPursuit} & \textbf{89.90} & 10.10 & 77.4 & \textbf{54/75$^\star$} \\
& \texttt{CEpsilonGreedy} & 88.08 & 11.92 & 79.4 & 21/75 \\
& \texttt{CEXP4} & 70.06 & 29.94 & 67.4 & 0/75 \\
& \texttt{CThompsonSampling} & 68.16 & 31.84 & 62.5 & 0/75 \\
& \texttt{CEpochGreedy} & 37.65 & 62.35 & 36.3 & 0/75 \\
& \cellcolor{findingBlue}\textit{Takeaway} & \multicolumn{4}{>{\raggedright\arraybackslash}p{0.48\textwidth}@{}}{\cellcolor{findingBlue}CPursuit is the strongest classical contextual baseline.} \\
\midrule
\multirow{6}{=}{\textbf{iCMABs}\newline internal corpus}
& \texttt{iCEpsilonGreedy} & \textbf{88.56} & 11.44 & 81.0 & \textbf{75/75$^\star$} \\
& \texttt{iCPursuit} & 68.69 & 31.31 & 56.9 & 0/75 \\
& \texttt{iCThompsonSampling} & 68.01 & 31.99 & 62.8 & 0/75 \\
& \texttt{iCEXP4} & 37.50 & 62.50 & 36.1 & 0/75 \\
& \texttt{iCEpochGreedy} & 37.57 & 62.43 & 36.1 & 0/75 \\
& \cellcolor{findingBlue}\textit{Takeaway} & \multicolumn{4}{>{\raggedright\arraybackslash}p{0.48\textwidth}@{}}{\cellcolor{findingBlue}iCEpsilonGreedy provides the strongest iCMAB floor and complete win dominance in this slice.} \\
\midrule
\multirow{4}{=}{\textbf{EXP3-based}\newline internal corpus}
& \texttt{GNeuralUCB} & \textbf{85.37} & 14.63 & 61.6 & \textbf{41/75$^\star$} \\
& \texttt{EXPNeuralUCB} & 80.86 & 19.14 & 18.0 & 28/75 \\
& \texttt{EXPUCB} & 78.06 & 21.94 & 68.8 & 6/75 \\
& \cellcolor{findingOrange}\textit{Takeaway} & \multicolumn{4}{>{\raggedright\arraybackslash}p{0.48\textwidth}@{}}{\cellcolor{findingOrange}GNeuralUCB is the strongest EXP-family representative, but its floor remains weaker than the top contextual variants.} \\
\midrule
\multirow{5}{=}{\textbf{Hybrid neural}\newline internal corpus}
& \texttt{iCPursuitNeuralUCB} & \textbf{90.86} & 9.14 & 76.4 & \textbf{23/75$^\star$} \\
& \texttt{CPursuitNeuralUCB} & 89.00 & 11.00 & 45.0 & 17/75 \\
& \texttt{GNeuralUCB} & 88.99 & 11.01 & 62.2 & 21/75 \\
& \texttt{EXPNeuralUCB} & 88.37 & 11.63 & 63.9 & 14/75 \\
& \cellcolor{findingGreen}\textit{Takeaway} & \multicolumn{4}{>{\raggedright\arraybackslash}p{0.48\textwidth}@{}}{\cellcolor{findingGreen}Hybrid pursuit-neural policies define the strongest internal performance ceiling.} \\
\midrule\midrule
\multirow{5}{=}{\textbf{\paperTwo}\newline external}
& \texttt{iCPursuitNeuralUCB} & \textbf{74.44} & 25.56 & 54.4 & \textbf{25/75$^\star$} \\
& \texttt{CPursuitNeuralUCB} & 73.15 & 26.85 & 48.3 & 8/75 \\
& \texttt{GNeuralUCB} & 73.17 & 26.83 & 46.7 & 20/75 \\
& \texttt{EXPNeuralUCB} & 71.29 & 28.71 & 41.6 & 22/75 \\
& \cellcolor{findingBlue}\textit{Takeaway} & \multicolumn{4}{>{\raggedright\arraybackslash}p{0.48\textwidth}@{}}{\cellcolor{findingBlue}iCPursuitNeuralUCB is the best average and win-count performer.} \\
\midrule
\multirow{5}{=}{\textbf{\paperSeven}\newline external}
& \texttt{iCPursuitNeuralUCB} & \textbf{77.50} & 22.50 & 28.7 & \textbf{58/75$^\star$} \\
& \texttt{GNeuralUCB} & 70.89 & 29.11 & 41.5 & 10/75 \\
& \texttt{CPursuitNeuralUCB} & 70.89 & 29.11 & 41.5 & 0/75 \\
& \texttt{EXPNeuralUCB} & 69.54 & 30.46 & 33.0 & 7/75 \\
& \cellcolor{findingGreen}\textit{Takeaway} & \multicolumn{4}{>{\raggedright\arraybackslash}p{0.48\textwidth}@{}}{\cellcolor{findingGreen}iCPursuitNeuralUCB dominates the dense 50-node external setting.} \\
\midrule
\multirow{5}{=}{\textbf{\paperTwelve}\newline external}
& \texttt{iCPursuitNeuralUCB} & \textbf{41.55} & 58.45 & 23.9 & \textbf{26/75$^\star$} \\
& \texttt{CPursuitNeuralUCB} & 41.19 & 58.81 & 22.6 & 13/75 \\
& \texttt{GNeuralUCB} & 41.13 & 58.87 & 22.8 & 14/75 \\
& \texttt{EXPNeuralUCB} & 40.18 & 59.82 & 21.7 & 22/75 \\
& \cellcolor{findingOrange}\textit{Takeaway} & \multicolumn{4}{>{\raggedright\arraybackslash}p{0.48\textwidth}@{}}{\cellcolor{findingOrange}All methods degrade under the larger fusion topology, but iCPursuitNeuralUCB remains best.} \\
\midrule
\multirow{5}{=}{\textbf{\paperEight}\newline external}
& \texttt{iCPursuitNeuralUCB} & \textbf{67.41} & 32.59 & \textbf{44.1} & 1/5 \\
& \texttt{EXPNeuralUCB} & 62.08 & 37.92 & 23.9 & \textbf{2/5$^\star$} \\
& \texttt{GNeuralUCB} & 61.81 & 38.19 & 36.0 & 0/5 \\
& \texttt{CPursuitNeuralUCB} & 61.65 & 38.35 & 35.2 & \textbf{2/5$^\star$} \\
& \cellcolor{findingBlue}\textit{Takeaway} & \multicolumn{4}{>{\raggedright\arraybackslash}p{0.48\textwidth}@{}}{\cellcolor{findingBlue}iCPursuitNeuralUCB leads average efficiency; EXP/CP pursuit variants split the small-slice win count.} \\
\bottomrule
\end{tabular}

\vspace{0.2cm}
{\scriptsize \textit{Note:} Internal corpus rows use the validated CMAB, iCMAB, EXP3, and Hybrid master datasets with \texttt{Default} allocator and reported run/scale settings. External rows use the standardized external-testbed corpus where available. Gap = 100 - Efficiency. Floor = worst-case efficiency. Wins denote experiment-level win counts; $^\star$ marks the row with the highest win count in its scope.}
\end{table*}
